\documentclass[11pt]{article}
\usepackage[margin=1in]{geometry}
\usepackage{amsmath,amssymb,amsthm,mathtools}
\usepackage{graphicx}
\usepackage[hidelinks]{hyperref}
\usepackage[T1]{fontenc}
\usepackage{lmodern}

\newtheorem{theorem}{Theorem}
\newtheorem{lemma}{Lemma}
\newcommand{\R}{\mathbb R}
\newcommand{\dom}{\operatorname{dom}}
\newcommand{\ran}{\operatorname{ran}}
\newcommand{\zer}{\operatorname{zer}}
\newcommand{\Id}{\operatorname{Id}}
\newcommand{\gra}{\operatorname{gra}}

\title{A finite-dimensional nonclosed zero set\\
for a sum of maximally monotone operators}
\author{Candidate full solution packet for arXiv:1110.4877}
\date{11 August 2026}

\begin{document}
\maketitle

\begin{center}
\fbox{\parbox{0.90\linewidth}{\textbf{Status.}
Candidate full solution, likely valid.  The construction gives an affirmative
answer to Remark~3.12 of Bauschke--Bo\c{t}--Hare--Moursi.}}
\end{center}

\section{Source question}

Bauschke, Bo\c{t}, Hare, and Moursi study the primal solution set
$Z=\zer(A+B)$ for maximally monotone operators on a real Hilbert space
\cite{source}.  Their infinite-dimensional example shows that $Z$ need not be
closed.  Remark~3.12 asks whether this can already happen in finite dimension
when $\dom A\cap\dom B\ne\varnothing$.

\paragraph{Transcription.}
``We do not know whether or not such counterexamples can reside in
finite-dimensional Hilbert spaces when $\dom A\cap\dom B\ne\varnothing$.''
The source further observes that at least one operator must be
nonparamonotone and that standard relative-interior sum qualifications must
fail.

\clearpage
\begin{figure}[!ht]
  \centering
  \includegraphics[width=\textwidth]{figures/open_problem_crop.png}
  \caption{Source-PDF page 10: the infinite-dimensional example and the
  finite-dimensional question in Remark~3.12.}
\end{figure}
\clearpage

\section{Definitions}

For a set-valued operator $T:\R^2\rightrightarrows\R^2$, write
$\gra T=\{(z,z^*):z^*\in Tz\}$.  It is \emph{monotone} when
\[
 \langle z-w,z^*-w^*\rangle\ge0
 \quad ((z,z^*),(w,w^*)\in\gra T),
\]
and \emph{maximally monotone} when its graph has no proper monotone
extension.  We use Minty's criterion: a monotone $T$ on a Hilbert space is
maximally monotone if and only if $\ran(\Id+T)$ is the whole space
\cite{minty}.

Let $H=\R\times(-\infty,0]$ and $L=\R\times\{0\}$.  Their normal cones use
the outward-normal convention; in particular,
\[
 N_H(x,y)=
 \begin{cases}
 \{(0,0)\},&y<0,\\
 \{0\}\times[0,\infty),&y=0,
 \end{cases}
 \qquad
 N_L(x,0)=\{0\}\times\R.
\]

\section{Proof intuition}

To make the zero set an open ray, let one operator live on the lower
half-plane but force its first coordinate to stay positive; let the second
operator be the normal cone to the horizontal axis.  The only difficult point
is maximality of an operator with nonclosed domain.  The vector field
\[
 F(x,y)=\left(\frac{y}{x^2},\frac1x\right)
\]
is designed for this: its off-diagonal derivatives are skew and cancel in the
symmetric Jacobian, while the singularity at $x=0$ prevents any missing graph
point.  More decisively, the resolvent equation reduces to a scalar continuous
function whose endpoint limits cover every target in $\R^2$.

\section{Counterexample}

\begin{theorem}
There exist maximally monotone operators
$A,B:\R^2\rightrightarrows\R^2$ with
$\dom A\cap\dom B\ne\varnothing$ such that
\[
 \zer(A+B)=(0,\infty)\times\{0\}.
\]
Consequently, the primal solution set of a sum of two maximally monotone
operators need not be closed even in a finite-dimensional Hilbert space.
\end{theorem}

\begin{proof}
Set $D=(0,\infty)\times(-\infty,0]$.  For $(x,y)\in D$, put
\[
 F(x,y)=\left(\frac{y}{x^2},\frac1x\right),
\]
and define
\[
 A(x,y)=F(x,y)+N_H(x,y)\quad ((x,y)\in D),
 \qquad A(x,y)=\varnothing\quad ((x,y)\notin D).
\]
Thus, explicitly,
\[
 A(x,y)=
 \begin{cases}
 \left\{(y/x^2,1/x)\right\},&x>0,\ y<0,\\
 \{0\}\times[1/x,\infty),&x>0,\ y=0,\\
 \varnothing,&\text{otherwise}.
 \end{cases}
\]

We first prove that $A$ is maximally monotone.  On the open half-plane
$\{x>0\}$,
\[
 DF(x,y)=
 \begin{pmatrix}
 -2y/x^3&1/x^2\\
 -1/x^2&0
 \end{pmatrix},
 \qquad
 \frac{DF+DF^{\mathsf T}}2=
 \begin{pmatrix}
 -2y/x^3&0\\0&0
 \end{pmatrix}.
\]
The latter matrix is positive semidefinite on $D$.  Integrating along the
line segment between any two points of the convex set $D$ shows that $F$ is
monotone on $D$.  The boundary term is a value of the monotone normal cone
$N_H$, so $A$ is monotone.

It remains to prove $\ran(\Id+A)=\R^2$.  Fix an arbitrary $(p,q)\in\R^2$.
If $q>0$ and $p\ge1/q$, take
\[
 (x,y)=(p,0),\qquad r=q-\frac1p\ge0.
\]
Then $(0,1/p+r)\in A(p,0)$ and
$(p,q)=(p,0)+(0,1/p+r)$.

In all other cases seek an interior point by setting
\[
 y=q-\frac1x<0.
\]
The first coordinate of $(x,y)+F(x,y)$ equals
\[
 h_q(x):=x+\frac{q}{x^2}-\frac1{x^3}.
\]
If $q>0$ and $p<1/q$, use $0<x<1/q$.  Then
\[
 \lim_{x\downarrow0}h_q(x)=-\infty,
 \qquad
 \lim_{x\uparrow1/q}h_q(x)=\frac1q>p.
\]
Continuity gives $x\in(0,1/q)$ with $h_q(x)=p$, and the definition gives
$y<0$.  If $q\le0$, use $x\in(0,\infty)$; now
\[
 \lim_{x\downarrow0}h_q(x)=-\infty,
 \qquad
 \lim_{x\to\infty}h_q(x)=+\infty,
\]
so again $h_q(x)=p$ for some $x>0$, while automatically
$y=q-1/x<0$.  In either interior case,
\[
 (x,y)+F(x,y)=(h_q(x),q)=(p,q).
\]
Hence $\ran(\Id+A)=\R^2$, and Minty's criterion proves that $A$ is maximally
monotone.

Now set $B=N_L$, which is maximally monotone.  The common domain is
\[
 \dom A\cap\dom B=(0,\infty)\times\{0\}.
\]
For $t>0$,
\[
 A(t,0)=\{0\}\times[1/t,\infty),
 \qquad B(t,0)=\{0\}\times\R,
\]
and therefore $(A+B)(t,0)=\{0\}\times\R$ contains $(0,0)$.  Outside the
common domain, $A+B$ is empty.  It follows that
\[
 \zer(A+B)=(0,\infty)\times\{0\},
\]
which is not closed in $\R^2$.
\end{proof}

\paragraph{Consistency with the source obstruction.}
The operator $A$ is not paramonotone.  Indeed, fix $x>0$ and distinct
$y_1,y_2<0$.  The monotonicity pairing between
$((x,y_i),F(x,y_i))$ is zero because the displacement is vertical while the
value difference is horizontal.  But $F(x,y_2)\notin A(x,y_1)$, so the
cross-membership required by paramonotonicity fails.

\section{Verification and novelty boundary}

The accompanying script samples graph pairs and arbitrary targets for the
Minty range equation; it is only a stress check, since the proof above is
exact.  On 11 August 2026, the run indexes and bounded current-literature
searches were queried using arXiv:1110.4877, the exact Remark~3.12 wording,
``$\zer(A+B)$ not closed'', ``maximally monotone'', and
``finite-dimensional''.  Nearby Attouch--Th\'era papers from 2013, 2016, and
2021 were checked at title/abstract level.  No later exact answer or occurrence
of this construction was found.  This supports, but cannot guarantee,
novelty.

\paragraph{Human review.}
The key check is the three-regime surjectivity argument for $\Id+A$, including
the endpoint $x=1/q$ when $q>0$.  The normal-cone sign at $y=0$ should then be
verified.  No computational or conditional dependency remains.

\begin{thebibliography}{9}
\bibitem{source}
H.~H. Bauschke, R.~I. Bo\c{t}, W.~L. Hare, and W.~M. Moursi,
\emph{Attouch--Th\'era duality revisited: paramonotonicity and operator
splitting}, arXiv:1110.4877 (2011), especially Remark~3.12.

\bibitem{minty}
G.~J. Minty,
\emph{Monotone (nonlinear) operators in Hilbert space},
Duke Math. J. \textbf{29} (1962), 341--346.
\end{thebibliography}

\end{document}
